\documentclass[11pt]{article}
\usepackage[margin=1in]{geometry}
\usepackage{amsmath,amssymb,amsthm}
\usepackage{hyperref}
\usepackage{graphicx}

\newtheorem{theorem}{Theorem}
\newtheorem{lemma}{Lemma}
\newtheorem{remark}{Remark}

\newcommand{\R}{\mathbb R}
\newcommand{\dd}{\mathrm d}
\newcommand{\dx}{\mathrm dx}
\newcommand{\mm}{\mathfrak m}
\newcommand{\Test}{\operatorname{Test}}
\newcommand{\TestForm}{\operatorname{TestForm}}
\newcommand{\Lip}{\operatorname{Lip}}
\newcommand{\RCD}{\operatorname{RCD}}
\newcommand{\Tr}{\operatorname{Tr}}

\title{A boundary counterexample to test-form density in \(W_C^{1,2}\)}
\author{Agent lane 14}
\date{2026-07-18}

\begin{document}
\maketitle

\section*{Source and question}

The source paper is Andrea Carbonaro, Luca Tamanini, and Dario Trevisan,
\emph{Boundedness of Riesz transforms on \(\RCD(K,\infty)\) spaces},
arXiv:2308.16294. In Section 2, after defining test functions and test
1-forms by
\[
  \TestForm(X)
  :=
  \left\{
    \sum_{i=1}^n g_i\,\dd f_i:
    n\in\mathbb N,\ f_i,g_i\in\Test(X)
  \right\},
\]
the source recalls the associated covariant and Hodge Sobolev spaces
\(W_C^{1,2}(T^*X)\) and \(W_H^{1,2}(T^*X)\). It then states that it is an
open problem whether \(\TestForm(X)\) is dense in these Sobolev spaces, and
denotes by \(H_C^{1,2}(T^*X)\), \(H_H^{1,2}(T^*X)\) the respective closures.

\begin{center}
\includegraphics[width=.94\linewidth]{figures/open_problem_crop.png}
\end{center}

\section*{Result}

The density statement fails for the covariant Sobolev space \(W_C^{1,2}\).
The obstruction is already present on the compact interval. Test functions
for the metric-measure interval satisfy Neumann boundary conditions, so their
differentials vanish at the endpoints. The maximal covariant Sobolev space
still contains the parallel 1-form \(\dx\), which has nonzero endpoint trace.

\begin{theorem}
Let \(I=[0,1]\) be equipped with the Euclidean distance and Lebesgue measure.
Then \(I\) is an \(\RCD(0,\infty)\) space and
\[
  \TestForm(I)
  \quad\hbox{is not dense in}\quad
  W_C^{1,2}(T^*I).
\]
More precisely, the \(W_C^{1,2}\)-closure of \(\TestForm(I)\) is contained in
\(H_0^1(I)\,\dx\), while the constant form \(\dx\) belongs to
\(W_C^{1,2}(T^*I)\).
\end{theorem}

\section*{Interval calculus}

We use the standard identification of first-order calculus on the interval:
\[
  W^{1,2}(I)=H^1(I), \qquad
  L^2(T^*I)=L^2(I)\,\dx.
\]
Under this identification, a 1-form is written \(u\,\dx\). The maximal
covariant Sobolev space is
\[
  W_C^{1,2}(T^*I)=H^1(I)\,\dx,
  \qquad
  \nabla(u\,\dx)=u'\,\dx\otimes\dx,
\]
and hence
\[
  \|u\,\dx\|_{W_C^{1,2}}^2
  =
  \int_0^1 |u|^2\,\dx+\int_0^1 |u'|^2\,\dx.
\]
This is the usual weak covariant derivative on the flat one-dimensional
manifold with boundary. It is also what the distributional definition of
\(\nabla\) gives on this metric-measure model, since test vector fields probe
the ordinary weak derivative in the interior and impose no vanishing condition
on the form being differentiated.

\begin{lemma}
If \(f\in D(\Delta)\) for the Cheeger Laplacian on \(I\), then
\(f\in H^2(0,1)\) and \(f'(0)=f'(1)=0\) in the trace sense.
\end{lemma}

\begin{proof}
The Cheeger energy on \(I\) is
\[
  \mathrm{Ch}(f)=\frac12\int_0^1 |f'|^2\,\dx
  \qquad (f\in H^1(I)).
\]
Thus \(f\in D(\Delta)\) means that there is \(h\in L^2(I)\) such that, up to
the source's harmless sign convention for \(\Delta\),
\[
  \int_0^1 h\varphi\,\dx=-\int_0^1 f'\varphi'\,\dx
  \qquad\forall \varphi\in H^1(I).
\]
Testing first with \(\varphi\in C_c^\infty(0,1)\) shows that the distributional
second derivative of \(f\) belongs to \(L^2(0,1)\), so \(f\in H^2(0,1)\). For
arbitrary \(\varphi\in H^1(I)\), integration by parts gives
\[
  -\int_0^1 f'\varphi'\,\dx
  =
  \int_0^1 f''\varphi\,\dx
  - f'(1)\varphi(1)+f'(0)\varphi(0).
\]
Comparing with the defining identity and using arbitrary endpoint values of
\(\varphi\) yields \(f'(0)=f'(1)=0\).
\end{proof}

\begin{lemma}
Every test 1-form on \(I\) has zero endpoint trace as an \(H^1\) coefficient.
That is, if \(\omega\in\TestForm(I)\) and \(\omega=u\,\dx\), then
\[
  u\in H^1(I),
  \qquad
  u(0)=u(1)=0.
\]
\end{lemma}

\begin{proof}
It is enough to consider one summand \(g\,\dd f\), with
\(f,g\in\Test(I)\). By the previous lemma, \(f'\in H^1(I)\) and
\(f'(0)=f'(1)=0\). Also \(g\in\Test(I)\) satisfies \(|\dd g|\in L^\infty\),
so on the interval \(g\) has a Lipschitz representative. The coefficient of
\(g\,\dd f\) is \(g f'\). Since \(g\in W^{1,\infty}(I)\) and \(f'\in H^1(I)\),
the product \(g f'\) belongs to \(H^1(I)\), and its endpoint traces are
\[
  (g f')(0)=g(0)f'(0)=0,
  \qquad
  (g f')(1)=g(1)f'(1)=0.
\]
Finite sums preserve the same property.
\end{proof}

\section*{Proof of the theorem}

Let
\[
  \mathcal C
  :=
  \overline{\TestForm(I)}^{\,W_C^{1,2}}
  \subset W_C^{1,2}(T^*I).
\]
Writing a form as \(u\,\dx\), the \(W_C^{1,2}\)-norm is exactly the \(H^1\)
norm of \(u\). The trace map
\[
  \Tr:H^1(I)\longrightarrow \R^2,
  \qquad
  \Tr(u)=(u(0),u(1)),
\]
is continuous. Since every test-form coefficient has trace \((0,0)\), every
element of the closure \(\mathcal C\) also has zero trace. Hence
\[
  \mathcal C\subset H_0^1(I)\,\dx.
\]

Now consider the constant 1-form
\[
  \omega_0=\dx.
\]
Its coefficient is \(u\equiv 1\), which belongs to \(H^1(I)\), and
\[
  \nabla\omega_0=0.
\]
Thus \(\omega_0\in W_C^{1,2}(T^*I)\). But
\[
  \Tr(1)=(1,1)\ne(0,0),
\]
so \(\omega_0\notin \mathcal C\). Therefore \(\TestForm(I)\) is not dense in
\(W_C^{1,2}(T^*I)\).

\section*{Verification notes}

The argument uses only the one-dimensional Sobolev trace theorem and the
standard identification of the Cheeger calculus on a compact Euclidean
interval. The key mechanism is the mismatch between:
\[
  \Test(I)\subset D(\Delta_{\mathrm{Neu}})
  \quad\Rightarrow\quad
  \dd f=f'\dx\hbox{ has zero endpoint trace},
\]
and
\[
  W_C^{1,2}(T^*I)=H^1(I)\dx,
\]
which contains forms with arbitrary endpoint trace.

\section*{Limitations and novelty check}

This packet addresses the covariant Sobolev space \(W_C^{1,2}\). It does not
settle the analogous density question for the Hodge Sobolev space
\(W_H^{1,2}\), where the codifferential may encode a different boundary
condition.

The claim should be read with the standard convention, also used in the source
passage, that \(W_C^{1,2}(T^*X)\) is the maximal space of \(L^2\) 1-forms whose
distributional covariant derivative lies in \(L^2\), while
\(H_C^{1,2}(T^*X)\) is the closure of test forms. If one redefined
\(W_C^{1,2}\) itself to mean the closure, the density statement would be
tautological; that is not the convention indicated by the source.

The bounded novelty check searched the run indexes for arXiv:2308.16294 and
for phrases around \(\TestForm\), density, \(W_C^{1,2}\), \(H_C^{1,2}\),
covariant Sobolev spaces, and \(\RCD\) spaces. Web searches for exact phrases
around ``TestForm dense'' and ``\(W_C^{1,2}\)'' found the source paper and
related background, but no later public answer or pre-existing interval
counterexample in this bounded search.

\section*{References}

A. Carbonaro, L. Tamanini, and D. Trevisan,
\emph{Boundedness of Riesz transforms on \(\RCD(K,\infty)\) spaces},
arXiv:2308.16294, 2023.

N. Gigli, \emph{Nonsmooth differential geometry -- an approach tailored for
spaces with Ricci curvature bounded from below}, Mem. Amer. Math. Soc. 251
(2018), no. 1196.

\end{document}
